% !TeX root = ../thesis.tex
\chapter{Introduction}
\label{chap:introduction}

\section{Motivation}
\label{sec:introduction-motivation}

Long, untrimmed surveillance video presents a difficult allocation problem: most
recorded time is routine, while events that warrant review may be brief, visually
ambiguous, and separated by long periods of normal activity. Automated video
analysis can help prioritize where a person looks, but an anomaly score alone is
not an emergency-response decision. A useful research account must distinguish
at least three questions: whether a video is anomalous, what kind of response a
predefined label scheme associates with it, and where the supporting event occurs
in time. Conflating these questions can make a strong detector appear to be a
capable response system even when its service selection or temporal grounding is
weak.

UCF-Crime made this setting accessible at scale by collecting 1,900 long,
untrimmed surveillance videos, including normal activity and 13 anomaly
categories, and by defining both binary anomaly detection and anomalous-activity
recognition tasks \cite{sultani2018real}. Its training supervision is available
at video level rather than as clip-level event labels. This reduces the need for
dense training annotation, but creates ambiguity: an anomalous video can contain
many normal clips, and its category does not identify which clips depict the
event. Weakly supervised video anomaly detection therefore has to learn from
bags of temporally mixed evidence rather than from precisely labeled instances.

The application framing raises a second difficulty. A camera record does not by
itself specify whether Police, Medical, Fire, some other service, or no immediate
service is appropriate. Such a decision depends on context that may not be
visible in the video and on institutional, legal, and professional judgment.
Research on semi-automated closed-circuit television (CCTV) has also shown that
system accuracy, confidence information, task complexity, and interface design
affect operator reliance and workload \cite{dadashi2013semiautomated}. For this
reason, response-oriented outputs should be evaluated as decision-support
signals with explicit human oversight, not presented as automatic dispatch.

The goal of this thesis is to build and evaluate a weakly supervised pipeline
that answers all three questions for UCF-Crime videos, and to measure each
answer separately against simple references. In the registered title,
\emph{understanding} therefore means three measurable outputs: a binary anomaly
decision, a set of project-defined response labels, and the temporal alignment
of attention with annotated events. \emph{Emergency response} names the
application framing of an offline system; it is not an operationally validated
capability.

\section{Problem Statement and Research Gap}
\label{sec:introduction-gap}

Weakly supervised video anomaly detection is commonly formulated with multiple
instance learning (MIL): a video is a bag, temporal snippets are instances, and
only the bag label is known during training \cite{sultani2018real,tian2021weakly}.
This formulation supports video-level learning and can produce snippet scores,
but success at one level does not establish success at another. A high video
area under the receiver operating characteristic curve (AUROC) does not
guarantee balanced recognition of anomaly categories, selective
multi-label response prediction, or accurate event boundaries. Likewise, a
visually plausible attention trace is not automatically a calibrated anomaly
probability or a faithful explanation of a decision.

Three methodological strands are relevant but remain conceptually distinct.
First, pretrained video representations can reduce the amount of task-specific
learning required; VideoMAE learns spatio-temporal representations through
masked video reconstruction \cite{tong2022videomae}. Low-rank adaptation (LoRA)
offers a parameter-efficient way to adapt selected projections while leaving the
pretrained weights frozen \cite{hu2021lora}. Second, MIL aggregation and
normality-oriented memories can prioritize instances when only bag labels are
available \cite{ilse2018attention,zhou2023ur}. Third, independent sigmoid outputs
and imbalance-aware objectives support non-exclusive labels
\cite{boutell2004learning,ridnik2021asymmetric}. Combining these elements in
software is not itself evidence that each element is beneficial, nor that their
outputs form a reliable integrated response system.

Published weakly supervised methods on UCF-Crime are compared mainly through
a single number, the frame-level AUC of an anomaly score
(Section~\ref{sec:background-comparison}). Three things are rarely reported
together with it: an output that states which services an event involves, a
per-video measurement of where the event is placed in time, and a comparison
of each learned output with a trivial rule on the same videos. Without such
references, a strong detector can look like a capable response system without
showing whether a rule that predicts every service after the same anomaly
decision would do as well. This thesis addresses that evaluation gap. It implements one sequential
pipeline that adapts a video encoder, freezes clip features, adds a
normal-prototype distance, trains an attention MIL model with a binary head and
three response channels, and renders an offline computer-aided dispatch
(CAD)-style record. It then evaluates the binary, response, and temporal outputs
separately, each against a matched reference, across repeated Stage-2 training.

\section{Research Questions}
\label{sec:introduction-rqs}

The following four questions define the empirical scope of the thesis:

\noindent\textbf{RQ1:} \emph{What video-level binary anomaly-detection
performance does the sequential LoRA–\allowbreak prototype–\allowbreak MIL pipeline achieve on
UCF-Crime under weak video-level supervision?}

\noindent\textbf{RQ2:} \emph{On UCF-Crime, what limitations are evident in the
reconciled historical 14-class predictions, and what error profile does the
project-defined multi-label emergency-response formulation exhibit?}

\noindent\textbf{RQ3:} \emph{How stable are binary and project-defined
response-channel predictions across Stage-2 training seeds under fixed feature
and evaluation artifacts on UCF-Crime, and how do matched response-label error
metrics differ with and without binary gating?}

\noindent\textbf{RQ4:} \emph{To what extent do the final checkpoints’ attention
outputs align with annotated anomalous intervals under the dense indexed-video
UCF-Crime evaluation protocol?}

RQ2 also covers an earlier 14-class formulation of the same pipeline, which is
reported as development history (Section~\ref{sec:results-historical}). RQ1
concerns the complete sequential configuration as executed, while its
reported variability is conditional on the frozen representation and prototype
artifacts. RQ2 describes two non-equivalent tasks separately and does not define
a before-and-after performance comparison. RQ3 treats gating as a matched
computational transformation of response-label predictions, not as evidence of
real dispatch effectiveness. RQ4 includes all three final checkpoints and uses
attention only as a temporal saliency signal whose overlap with annotations must
be measured.

\section{Scope and Contributions}
\label{sec:introduction-contributions}

The work makes an implementation-and-evaluation contribution. It does not claim
to invent VideoMAE, LoRA, MIL, attention pooling, asymmetric loss, or CAD. Its
contributions are as follows:

\begin{enumerate}
    \item \textbf{A sequential weakly supervised implementation.} The pipeline
    links video indexing and clip preparation, LoRA adaptation of VideoMAE,
    frozen adapted features, a normal-prototype distance, guided bag
    construction, and a shared-attention MIL model with a binary head and three
    response channels. It also supports dense indexed evaluation and
    deterministic offline CAD-style output. The stages are trained one after
    another, not jointly from pixels.

    \item \textbf{A separated evaluation design.} Historical 14-class outputs,
    validation-only studies, fixed-bag test results, and dense indexed results
    are reported separately. Fixed manifests, three Stage-2 seeds, per-video
    records, matched populations, trivial baselines, and hash-identified local
    artifacts keep positive and negative findings traceable.

    \item \textbf{A matched characterization of system behavior.} The response
    head is compared with predicting every service, with and without the same
    binary gate. The attention signal is compared with whole-video, uniform,
    random, and prototype-distance references, and alignment is measured per
    video under near-complete indexed coverage. Two validation-only studies, one
    on the prototype score with guided sampling and one on normal bags in the
    response loss, support the reading of these comparisons.

    \item \textbf{Explicit operational boundaries.} The thesis separates
    derived response labels from observed emergency decisions, temporal saliency
    from explanation, source-video timestamps from compute latency, and an
    offline record from a live CAD connection. These limits define what the
    evidence can support.
\end{enumerate}

The scope follows from these choices. UCF-Crime is used because it combines
real surveillance footage, named incident categories, weak training labels, and
annotated test intervals (Section~\ref{sec:background-tasks}). The Police,
Medical, and Fire targets are deterministic project transformations of its
categories, not observed service assignments. Only Stage-2 training is repeated
across three seeds; the adapted feature store and prototype memory remain fixed.
Temporal evaluation measures alignment with annotated intervals, not causal
explanation or real-time localization. All conclusions are restricted to
UCF-Crime, the fixed splits and artifacts, and the executed protocols; the work
does not establish generalization to other sites, calibration, real-time
throughput, human-factors performance, safety, or deployment readiness.

\section{Thesis Structure}
\label{sec:introduction-structure}

Chapter~\ref{chap:background} develops the concepts needed to read the method
and evidence: UCF-Crime supervision, task boundaries, MIL, video representation
learning, LoRA, prototype guidance, attention pooling, multi-label prediction,
gating, and temporal metrics, and closes with a comparison of related
approaches. Chapter~\ref{chap:methods}
specifies the implemented pipeline, data identities, objectives, evaluation
paths, and reproducibility limits. Chapter~\ref{chap:results} reports the
historical, fixed-bag, validation-only, dense indexed, and temporal results
without extending them into deployment claims. Chapter~\ref{chap:discussion}
answers RQ1--RQ4, interprets the error profiles, and addresses limitations,
ethics, human oversight, and future work. Chapter~\ref{chap:conclusion}
provides the final synthesis of the evidence and its bounded contribution.
